%% The eddy diffusion coefficient K_zz: what it is, how it is derived, how it
%% is measured, and what any of that fixes for the LHS 1140 b wind base that
%% EXHALE's binary H/He element-diffusion operator sits on.
%% Companion documents: LHS1140b/kzz_literature.md (the survey and the
%% verification status of every published number), LHS1140b/kzz_decision.md
%% (the D_eff structure, the He_Kzz scan and the adopted value),
%% docs/binary_diffusion_design.md (the operator itself).
%% Author: Kwang-Il Seon
\documentclass[english,a4paper]{my_memo}
\synctex=1
\usepackage{booktabs}
\usepackage{amsmath}
\usepackage{graphicx}
\usepackage{url}
\usepackage{babel}
\urlstyle{tt}
\emergencystretch=3em

\title{The eddy diffusion coefficient $K_{zz}$:\\
       theory, empirical determinations, and what fixes it\\
       at the LHS\,1140\,b wind base}
\author{Kwang-Il Seon}
\date{Last updated: \docmoddate}

\begin{document}
\maketitle

\section{Why this memo exists, and the answer up front}

EXHALE's binary H/He element-diffusion operator
(\path{src/modules/functions/binary_element_diffusion.f90},
formulated in \path{docs/binary_diffusion_design.md}) carries one number
that the code cannot derive from anything it solves: the eddy diffusion
coefficient \texttt{He\_Kzz}. On LHS\,1140\,b that number moves the
He\,\textsc{i}\,$10830$ equivalent width by four orders of magnitude and the
escape rate by $0.23$\,dex (\path{LHS1140b/kzz_decision.md}, section~3),
so it is not a detail. This memo collects what the physics, the modelling
literature and the measurements actually say about $K_{zz}$, and asks what
of it survives contact with a $226$\,K, $1\,\mu$bar base on a temperate
super-Earth.

\medskip
\noindent\textbf{Summary.} $K_{zz}$ is a one-dimensional closure, not a
material property. It is defined by the ratio of a horizontally averaged
tracer flux to a horizontally averaged mixing-ratio gradient, so it is a
property of the pair (atmosphere, tracer) and not of the atmosphere alone;
Zhang \& Showman (2018a,b) show that different species in the same
atmosphere have different $K_{zz}$ profiles and that the inferred value can
even be negative when transport is not diffusive. Three theoretical
constructions dominate: a convective mixing-length branch
($K = w_* L$, deep and irrelevant at $1\,\mu$bar), a breaking-gravity-wave
branch that supplies the familiar $K \propto P^{-1/2}$ growth
(Lindzen 1981), and effective coefficients read off three-dimensional
tracer transport in GCMs (Parmentier et al.\ 2013; Charnay et al.\ 2015).
Empirical determinations come from solar-system homopause occultations
(Festou et al.\ 1981: $1.4\times10^{6}$\,cm$^{2}$\,s$^{-1}$ at Jupiter's
homopause), from chemical quenching (Prinn \& Barshay 1977, the founding
case), and, on exoplanets, from cloud and haze lofting arguments that give
lower bounds only.

Two conclusions matter for the present application, and both are negative.
First, \emph{every} published number that can be evaluated at $1\,\mu$bar
comes from an atmosphere warmer and more strongly forced than this one, and
the only scaling that has ever been evaluated with LHS\,1140\,b's own
gravity and temperature (Arfaux \& Lavvas 2023) lands one to two decades
below the hot-Jupiter practice the code currently inherits. Second, and
less obvious, a \emph{constant} $K_{zz}$ is the wrong shape of assumption
for this problem regardless of its magnitude: the molecular coefficient it
competes against rises by nearly five decades between the base and
$1.2\,R_p$, so no single constant can both mix the base and follow the
molecular coefficient outward, while Lindzen's own result is that the eddy
coefficient should \emph{turn over} rather than keep growing. The adopted
$\texttt{He\_Kzz} = 10^{9}$\,cm$^{2}$\,s$^{-1}$ is defensible only as a
stated convention inherited from hot-Jupiter work, and it should be
reported alongside the $10^{8}$--$10^{10}$ bracket, not instead of it.

\medskip
\noindent This memo does not repeat the survey. The verification status of
every published number quoted here is in
\path{LHS1140b/kzz_literature.md}; the $D_{\rm eff}$ structure, the
\texttt{He\_Kzz} scan and the adopted value are in
\path{LHS1140b/kzz_decision.md}; the operator and its
$(D+K)$ form are in \path{docs/binary_diffusion_design.md},
sections~2.3 and~9.4. What is new here is the physics behind the numbers,
four newly verified primary citations (Sect.~\ref{sec:verif}), and an
assessment of what a principled determination would require.

\section{What $K_{zz}$ is}
\label{sec:what}

\subsection{The closure}

Take a three-dimensional atmosphere, a tracer with mixing ratio $q$, and
average horizontally over a level surface. The horizontally averaged
vertical tracer flux contains a mean-flow part and an eddy correlation
$\overline{\rho\,w'q'}$. The eddy diffusion closure asserts that the second
is proportional to the mean gradient,
%
\begin{equation}
  \overline{\rho\,w'q'} \;=\; -\,\bar\rho\,K_{zz}\,\frac{d\bar q}{dz},
  \qquad\text{i.e.}\qquad
  K_{zz} \equiv -\,\frac{\overline{\rho\,w'q'}}{\bar\rho\,d\bar q/dz},
  \label{eq:def}
\end{equation}
%
which is a \emph{definition} of $K_{zz}$ whenever both sides are measurable,
and a \emph{model} only when $K_{zz}$ is asserted in advance. Everything in
Sect.~\ref{sec:theory} is an attempt to predict the right-hand side without
running the three-dimensional flow.

Two features of Eq.~(\ref{eq:def}) are worth stating explicitly because
they are what separate eddy from molecular diffusion.

\paragraph{The gradient is in mixing ratio, not density.} A uniformly mixed
atmosphere has $d\bar q/dz = 0$ and therefore zero eddy flux, however steep
its density gradient. Molecular diffusion, by contrast, drives each species
toward its \emph{own} scale height and so has a non-zero flux in a uniformly
mixed atmosphere. Eddy mixing transports the mixture as a whole; molecular
diffusion separates it. This is why the aeronomical transport equation
(Banks \& Kockarts 1973) carries $D+K_{zz}$ on the gradient term and $D$
alone on the settling term, which is exactly the form EXHALE implements
(\texttt{binary\_diffusion\_design.md}, eq.~4):
%
\begin{equation}
  J_{\rm total} \;=\; J_{\rm mol} \;-\; \rho\,K_{zz}\,\frac{dX}{dr},
  \label{eq:DplusK}
\end{equation}
%
with $X$ the helium mass fraction and $J_{\rm mol}$ the binary molecular
flux of the design memo's eq.~(2)--(3).

\paragraph{The homopause is where the two are equal.} Setting
$K_{zz} = D$ defines the homopause. Below it the composition is held at the
reservoir value; above it the elements separate on their own scale heights.
The homopause is therefore not a property of $K_{zz}$ alone but of the
crossing of two profiles, and where a model places it is the single most
consequential thing $K_{zz}$ does. In the LHS\,1140\,b runs the homopause
radius moves from below the base at $K_{zz} = 0$ to $1.44\,R_p$ at
$K_{zz} = 10^{11}$ (\texttt{kzz\_decision.md}, section~3).

\subsection{$K_{zz}$ is a property of the (atmosphere, tracer) pair}
\label{sec:zs}

Equation~(\ref{eq:def}) has a tracer in it. Zhang \& Showman (2018a,b)
develop this into the central theoretical caveat of the subject: the
effective $K_{zz}$ inferred from a tracer depends on that tracer's chemical
lifetime, on the strength of the large-scale circulation and on horizontal
eddy mixing, so that ``different chemical species in a single atmosphere
should in principle have different eddy diffusion profiles'', and, where the
transport is not diffusive at all --- a mean overturning circulation, for
instance, whose upwelling and downwelling branches carry different tracer
abundances --- the $K_{zz}$ that reproduces the mean profile can be
\emph{negative}.

The consequence for us is specific and is not addressed anywhere in the
literature surveyed. The tracer EXHALE diffuses is the elemental He/H ratio:
chemically inert, of unbounded lifetime, and subject to a settling velocity
of its own. A $K_{zz}$ calibrated on CO/CH$_4$ quench chemistry, or on a
settling condensate, is calibrated on a tracer with a finite lifetime; in
Zhang \& Showman's framework a longer-lived tracer generally samples the
circulation more completely and yields a \emph{different} effective
coefficient. Whether the elemental ratio should take a larger or smaller
$K_{zz}$ than the quench tracers is not something this memo can settle;
it is recorded as a systematic that the transfer of literature values into
Eq.~(\ref{eq:DplusK}) silently assumes away.

\subsection{Definitional slack}

Three further ambiguities affect numbers taken from print:
%
\begin{itemize}
\item \emph{Which mixing ratio.} Eq.~(\ref{eq:def}) written for a molar
      mixing ratio and for a mass fraction differ when the mean molecular
      weight varies with height --- exactly the situation above a homopause.
      EXHALE's eddy term mixes toward uniform \emph{mass} fraction
      (\texttt{parameters.f90}, the \texttt{he\_kzz} comment).
\item \emph{Units.} $1$\,m$^{2}$\,s$^{-1} = 10^{4}$\,cm$^{2}$\,s$^{-1}$, and
      the literature uses both. \texttt{kzz\_literature.md}, section~2.1
      documents a place where the two appear to have been confused in print.
\item \emph{Where it is quoted.} A value is meaningless without the pressure
      at which it applies, unless the source asserts a plateau. Half the
      disagreement in Table~\ref{tab:lit} is a disagreement about height,
      not about magnitude.
\end{itemize}

\section{Theory}
\label{sec:theory}

\subsection{Convective mixing length}
\label{sec:ml}

In a convecting layer the natural estimate is $K_{zz} \simeq w_* L$ with
$w_*$ a convective velocity scale and $L$ a mixing length. Writing $w_*$
in terms of the convective heat flux gives the free-convection form of
Gierasch \& Conrath (1985), which Ackerman \& Marley (2001) adopt as their
eq.~(5):
%
\begin{equation}
  K \;=\; \frac{H}{3}\left(\frac{L}{H}\right)^{4/3}
          \left(\frac{R\,F}{\mu\,\rho_a\,c_p}\right)^{1/3},
  \qquad w_* = K/L,
  \label{eq:gc}
\end{equation}
%
with $H = RT/\mu g$ the scale height, $F$ the convective heat flux, $\rho_a$
the density and $c_p$ the specific heat. They state that the coefficient
$K$ ``is assumed to be the same as that for heat as derived for free
convection'', that ``the mixing length is typically assumed to be the
pressure scale height'' in freely convecting regions, and that $w_*$ is
``the convective velocity scale from mixing-length theory''. The same
expression is the one Charnay et al.\ (2015) write as their eq.~(16).

Three properties of Eq.~(\ref{eq:gc}) carry into everything downstream.

\paragraph{It is weakly sensitive to the energetics and strongly sensitive
to the scale height.} $K \propto F^{1/3}$ and $K \propto H$ at fixed $L/H$.
A factor $10^{3}$ in heat flux buys one decade in $K$; a factor $10$ in
scale height buys one decade as well. For a cold, high-gravity,
helium-dominated atmosphere both push the same way and the scale height
pushes harder.

\paragraph{Its normalization is admitted to be free.} Ackerman \& Marley
adopt the prefactor $1/3$ ``based on previous modeling studies of the Jovian
atmosphere'' and state plainly that ``the constant coefficient scaling the
eddy diffusion coefficient is only loosely constrained by observations''.
Their baseline Jovian values just below the ammonia cloud are $H = L =
20$\,km, $K = 2\times10^{8}$\,cm$^{2}$\,s$^{-1}$ and $w_* =
1.1$\,m\,s$^{-1}$; they also impose a floor $K_{\rm min} =
10^{5}$\,cm$^{2}$\,s$^{-1}$, writing that ``we also assume an eddy diffusion
coefficient no less than a prescribed minimum value'', which the paper
attributes to ``residual turbulence'' from ``breaking buoyancy waves'',
citing Lindzen (1981). (Those numerals and the attribution are read from the
arXiv version; the sentences quoted are confirmed present in the published
text, and the two fragments in the last clause are individually confirmed
--- see Sect.~\ref{sec:verif}.)

\paragraph{It is a deep-atmosphere statement, and the flux in it is the
internal one.} Convection carries the interior heat flux; the absorbed
stellar flux is re-radiated from the photosphere and does not pass through
the convective zone. Ackerman \& Marley evaluate $F = \sigma T_{\rm
eff}^{4}$ with Jupiter's \emph{internal} $T_{\rm eff} = 124$\,K. Charnay et
al.\ (2015) warn that this formula ``strongly overestimates the convective
heat flux for irradiated planets'' when $F_c$ is set from $T_{\rm eff}$.
This corrects a loose step in \path{LHS1140b/kzz_literature.md},
section~4, which compares LHS\,1140\,b and HD\,209458\,b through their
\emph{insolation} on the convective branch: the insolation ratio bears on
the wave-driving branch of Sect.~\ref{sec:gw}, through the day/night
contrast, and not on Eq.~(\ref{eq:gc}). The direction of the conclusion in
that document is unaffected --- both branches point the same way --- but the
convective comparison should be made through the internal heat flux, which
is unconstrained for LHS\,1140\,b and presumably small for an old,
low-mass, temperate planet.

\paragraph{And it does not apply at the wind base.} The LHS\,1140\,b base is
at $1\,\mu$bar in a stably stratified, radiatively controlled layer, some
$80$--$90$ isothermal scale heights above the $1$\,bar level
($H = kT/\mu m_H g = 25$\,km at $T = 226$\,K, $\mu = 4$, $g =
1837$\,cm\,s$^{-2}$; $H/R_p = 2.3\times10^{-3}$). Equation~(\ref{eq:gc})
sets the \emph{deep} value on which the wave branch of the next section is
built; it says nothing directly about the base.

\subsection{Breaking gravity waves and the $P^{-1/2}$ growth}
\label{sec:gw}

The mechanism that supplies mixing in a stably stratified upper atmosphere
is wave breaking, and its canonical treatment is Lindzen (1981). The
argument in outline: an internal gravity wave propagating upward conserves
its vertical flux of wave energy, so its velocity amplitude grows as
$\rho^{-1/2}$, i.e.\ roughly as $P^{-1/2}$; the growth cannot continue
indefinitely, because at large enough amplitude the wave-induced shear or
lapse rate makes the flow unstable and the wave breaks; the turbulence
produced by breaking is then parameterized as an eddy coefficient of exactly
the size needed to hold the wave at its saturation amplitude.

Lindzen's abstract states the premise in those terms --- the expectation
that ``sufficient turbulence would be generated to prevent the growth of
wave amplitude with height (roughly as (basic pressure)$^{-1/2}$)'' --- and
the paper's stated purpose is to extend that picture to smaller-amplitude
waves, to the effect of mean winds on the waves, and to the waves' effect
on the mean momentum budget.

Two things follow, and the second is usually dropped.

\paragraph{The slope.} A saturated wave spectrum gives an eddy coefficient
that \emph{grows} with height, and the $\rho^{-1/2}$ amplitude growth is
where the exponent comes from. Every power-law $K_{zz}(P)$ in the exoplanet
literature is of this family: Parmentier et al.\ (2013) fit
$K_{zz} = 5\times10^{4}\,P_{\rm bar}^{-1/2}$\,m$^{2}$\,s$^{-1}$ to their
HD\,209458\,b GCM tracers; Charnay et al.\ (2015) fit
$K_{zz} = K_{zz,0}\,P_{\rm bar}^{-0.4}$\,m$^{2}$\,s$^{-1}$ on GJ\,1214\,b
and justify the exponent by noting that planetary waves ``generally grow as
$P^{-1/2}$'', expecting ``an eddy diffusion coefficient with an exponent
between $-1/3$ and $-1/2$''.

\paragraph{The turnover.} Lindzen's own conclusion is that the eddy
coefficient does \emph{not} rise without limit. The mean wind filters the
wave spectrum before it arrives, critical levels remove waves, and the paper
reports that ``eddy coefficients may peak'' at particular altitudes and
then ``decrease with height rather sharply above these levels'' (terrestrial
mesosphere; the abstract quotes $50$\,km in winter mid-latitudes and
$70$--$80$\,km in summer). Extrapolating a saturated-wave power law upward
past the level where the wave field is exhausted is therefore not supported
by the theory that produced the power law. This is the physical content of
the plateau assumption used by Moses et al.\ (2011) --- constant $K_{zz}$
above the GCM top --- and of the homopause plateau of Arfaux \& Lavvas
(2023) in Sect.~\ref{sec:homo}, and it is the reason the extrapolation of
Charnay's $P^{-0.4}$ fit $1.5$ decades above its GCM top
(\texttt{kzz\_literature.md}, section~3) should be read as an upper bound.

\paragraph{And the magnitude is set by the source, not by the slope.}
Wave breaking fixes how $K_{zz}$ varies with pressure; what fixes its value
is how much wave flux is launched from below. On tidally locked, strongly
irradiated planets that source is the day/night thermal contrast --- which
is what Parmentier et al.\ (2013) identify, finding that ``vertical mixing
results not from small-scale convection but from the large-scale circulation
driven by the day-night heating contrast''. LHS\,1140\,b's day/night
contrast, in absolute terms, is smaller than a hot Jupiter's by orders of
magnitude, so the same mechanism should be correspondingly weaker even if
the slope survives. This is the branch on which the insolation comparison
of \texttt{kzz\_literature.md}, section~4 is the relevant one.

\subsection{Effective diffusion measured in a GCM}
\label{sec:gcm}

The most defensible route to $K_{zz}$ for a given planet is to stop
predicting it and to measure it in a three-dimensional model, through
Eq.~(\ref{eq:def}). The methodology, as used by Parmentier et al.\ (2013)
and Charnay et al.\ (2015), is:

\begin{enumerate}
\item Run a GCM to a statistically steady circulation.
\item Advect one or more \emph{passive} tracers in that flow. Passive means
      the tracer does not feed back on the dynamics; it may still carry a
      source/sink or, in Charnay's case, a settling velocity representing
      condensate particles of a given size.
\item Average the tracer field horizontally at each level to get $\bar q(z)$
      and the eddy flux $\overline{\rho w'q'}(z)$.
\item Form the flux/gradient ratio of Eq.~(\ref{eq:def}), or equivalently
      fit a one-dimensional diffusion model to reproduce $\bar q(z)$, and
      report the result as a power law in pressure.
\end{enumerate}

Two results of that exercise are more important than the fitted constants.

First, the effective coefficient is much smaller than the naive estimate.
Parmentier et al.\ report their GCM-derived $K_{zz}$ to be ``two orders of
magnitude smaller than what is obtained when multiplying the vertical scale
height by the root mean square of the vertical velocity'' --- the
$w_{\rm rms}H$ estimate. The reason is that $w_{\rm rms}H$ counts all
vertical motion as mixing, including the reversible up-and-down of a
large-scale circulation that returns the tracer to where it came from.
The flux/gradient ratio counts only the part that is correlated with the
tracer anomaly.

Second, and following from it, the tracer matters (Sect.~\ref{sec:zs}):
Charnay's coefficients depend on the settling velocity of the particle the
tracer represents, and are quoted as separate fits for different atmospheric
compositions.

Neither GCM covers a $1\,\mu$bar level. Charnay's model top is $3$\,Pa
($3\times10^{-5}$\,bar); Parmentier's parameterization is stated valid
``between $\sim 1$\,bar and $\sim 1$\,microbar'', which is the only entry in
the survey whose stated validity reaches the pressure of interest.

\subsection{Homopause scaling}
\label{sec:homo}

The construction closest in spirit to what a wind model needs is a
prediction of the coefficient \emph{at the homopause}, taken constant above
it. Arfaux \& Lavvas (2023) build a three-branch parameterization --- a deep
convective branch, a gravity-wave middle branch, and a constant homopause
plateau --- and write the plateau value as a scaling from Jupiter:
%
\begin{equation}
  K_{\rm top} \;=\; K_{zz}^{J}\,\frac{v}{v_J}\,\frac{g_J}{g}\,\frac{T}{T_J},
  \qquad K_{zz}^{J} = 10^{7}\ {\rm cm^{2}\,s^{-1}} ,
  \label{eq:al}
\end{equation}
%
with $v$ a characteristic turbulent velocity ($v/v_J = 5$ retained for
HD\,209458\,b). The $T/g$ dependence is transparent once read as
$K \sim vH$: the mixing length at the homopause is the scale height,
$H = kT/\mu m_H g$, so at fixed composition $K_{\rm top} \propto vT/g$
directly. (That reading is mine, not a statement quoted from the paper.)

Evaluated with LHS\,1140\,b's own $g = 18.4$\,m\,s$^{-2}$ and $T = 226$\,K
against Jupiter's $g_J = 24.79$\,m\,s$^{-2}$, $T_J = 110$\,K
(\texttt{kzz\_literature.md}, section~3), Eq.~(\ref{eq:al}) gives
$2.8\times10^{7}$\,cm$^{2}$\,s$^{-1}$ for $v/v_J = 1$ and
$1.4\times10^{8}$ for $v/v_J = 5$. The same formula applied to
HD\,209458\,b ($g = 9.4$\,m\,s$^{-2}$, $T = 1450$\,K, $v/v_J = 5$) returns
$1.7\times10^{9}$, within a factor $2$ of the $10^{9}$\,cm$^{2}$\,s$^{-1}$
that Taylor et al.\ (2025) adopt on the strength of this parameterization,
so the scaling reproduces the hot-Jupiter practice it is meant to reproduce
and its much lower LHS\,1140\,b prediction is not an artifact of
misapplication.

Its weak point is the anchor. Section~\ref{sec:homo-obs} shows that the
measured Jovian homopause coefficient is itself uncertain by nearly an
order of magnitude, and that uncertainty propagates through
Eq.~(\ref{eq:al}) unchanged.

\section{Empirical determinations}
\label{sec:emp}

\subsection{Solar-system homopause occultations}
\label{sec:homo-obs}

The only direct measurements of $K_{zz}$ at a homopause are solar-system
occultations. Festou et al.\ (1981) is the standard Jovian case: Voyager~2
watched the star Regulus set behind Jupiter's atmosphere on 9 July 1979, and
the absorption in the $910$--$1200$\,\AA{} range (H$_2$ Lyman and Werner
bands) and in the $1250$--$1600$\,\AA{} range (hydrocarbons) gives density
profiles for H$_2$, CH$_4$ and C$_2$H$_6$ directly. The heavier species fall
off faster than H$_2$ above the homopause, and where the transition occurs
fixes $K_{zz}$. Their published abstract states that the analysis
``yields a value of the eddy diffusion coefficient at the homopause'' of
%
\begin{equation}
  K_{\rm homopause}({\rm Jupiter,\ equatorial})
  = 1.4^{+0.8}_{-0.7}\times10^{6}\ {\rm cm^{2}\,s^{-1}},
\end{equation}
%
consistent, they note, with the value deduced from the H\,Ly$\alpha$ and
He\,$584$\,\AA{} emission data.

The method is the cleanest available: no chemical network, no free
efficiency parameter, and the same physics (a species-mass-dependent scale
height above a mixing level) that Eq.~(\ref{eq:DplusK}) implements. Its
limitations are equally clear: it measures $K_{zz}$ at one level, on one
planet, in one region --- the paper is explicit that the value is for the
equatorial region --- and it presupposes that a diffusive description is
appropriate, which is the assumption Sect.~\ref{sec:zs} questions.

\medskip
\noindent\textbf{This number matters here for a reason beyond its own
interest.} Arfaux \& Lavvas anchor Eq.~(\ref{eq:al}) at
$K_{zz}^{J} = 10^{7}$\,cm$^{2}$\,s$^{-1}$, after Koskinen et al.\ (2010) as
recorded in \texttt{kzz\_literature.md} (that attribution is not verified
here). That is a factor $\simeq 7$ above Festou et al.'s measured value.
The exoplanet homopause scaling therefore inherits a factor of several
before any planet-to-planet extrapolation begins. If the anchor were taken
at the Festou value instead, Eq.~(\ref{eq:al}) would give
$4\times10^{6}$--$2\times10^{7}$\,cm$^{2}$\,s$^{-1}$ for LHS\,1140\,b, close
to the molecular coefficient at the base (Sect.~\ref{sec:base}) --- i.e.\
close to no eddy mixing at all in this problem. This is stated as a
sensitivity, not as a recommendation: the two Jovian numbers may refer to
different levels, regions or epochs, and reconciling them is outside this
memo.

\subsection{Chemical quenching}
\label{sec:quench}

The second empirical route is disequilibrium chemistry. A species whose
chemical relaxation time $t_{\rm chem}$ is short compared with the mixing
time $t_{\rm mix} \simeq H^{2}/K_{zz}$ tracks local chemical equilibrium; a
species for which the inequality reverses is frozen, or ``quenched'', at the
abundance it had where $t_{\rm chem} = t_{\rm mix}$. The observed abundance
therefore measures $K_{zz}$ at the quench level, provided the chemical
network is known.

Prinn \& Barshay (1977) is the founding application, on CO in Jupiter's
visible atmosphere. Their published abstract states that CO's presence
``is a direct result of very rapid upward mixing from levels in the deep
atmosphere where the temperature is about $1100$ degrees K and where carbon
monoxide is thermodynamically much more stable'', that ``as a consequence
the observed carbon monoxide mixing ratio is a sensitive function of the
vertical eddy mixing coefficient'', and that the inferred coefficient is
``about three to four orders of magnitude greater than that in the earth's
troposphere''.

The method is now standard on exoplanets: Moses et al.\ (2011) use it,
together with mixing-length estimates from a GCM, to build the $K_{zz}$
profiles for HD\,189733\,b and HD\,209458\,b that Koskinen et al.\ (2013b)
and much subsequent work quote.

Its two limitations bear directly on the present problem. The quench level
sits at bars to tens of millibars --- deep --- and the value obtained there
has to be extrapolated by many decades to reach $1\,\mu$bar, which is where
the plateau assumptions of Moses et al.\ and of Sect.~\ref{sec:homo} enter.
And the tracer has, by construction, a finite chemical lifetime comparable
to the mixing time, which is precisely the regime in which Zhang \& Showman
(2018a,b) find the inferred coefficient to be tracer-dependent
(Sect.~\ref{sec:zs}).

\subsection{Cloud and haze lofting}

A third empirical handle gives bounds rather than values. Particles of a
given size settle at a computable terminal velocity; if they are observed
aloft, $K_{zz}$ must be large enough to keep them there. Ackerman \& Marley
(2001) turn this into a cloud model by balancing turbulent diffusion against
sedimentation with an adjustable sedimentation efficiency $f_{\rm rain}$;
Koskinen et al.\ (2013b) use the same reasoning in the other direction on
HD\,209458\,b, arguing that $K_{zz} \gtrsim 10^{9}$\,cm$^{2}$\,s$^{-1}$ is
required to prevent cloud settling through the cold trap. The output of this
route is a one-sided constraint entangled with an unknown particle size, and
it lands at millibar levels, not at $1\,\mu$bar.

\subsection{Can He\,10830 itself constrain $K_{zz}$?}
\label{sec:he-constraint}

There is a fourth possibility, specific to this problem: the He\,$10830$
line is made from helium that has crossed the homopause, so its strength
measures how much helium got through. The scan of
\path{LHS1140b/kzz_decision.md}, section~3 is exactly that experiment,
and it establishes both the sensitivity and its limit.

\begin{table}[htbp]\centering
\caption{Response of the LHS\,1140\,b model to \texttt{He\_Kzz}, from
\protect\path{LHS1140b/kzz_decision.md}, section~3. All rows are JFNK steady
states at He/H\,$=0.55$ with \texttt{He\_diffusion: True}, identical in
every other respect; the homopause radius is where the run's own
$D_{\rm eff}$ equals its $K_{zz}$. The last two rows are the
diffusion-off reference at the same reservoir composition and the
measurement of Cherubim et al.\ (2026).}
\label{tab:scan}
\begin{tabular}{lrrrr}
\toprule
$K_{zz}$ [cm$^{2}$\,s$^{-1}$] & homopause [$R_p$] &
$\log_{10}\dot M$ [g\,s$^{-1}$] & red depth [\%] & red EW [\%\,\AA] \\
\midrule
$0$        & below the base & $7.61$ & $9.97\times10^{-5}$ & $2.5\times10^{-5}$ \\
$10^{6}$   & $1.0004$ & $7.68$ & $1.00\times10^{-2}$ & $2.27\times10^{-3}$ \\
$10^{7}$   & $1.0108$ & $7.68$ & $2.57\times10^{-2}$ & $5.78\times10^{-3}$ \\
$10^{8}$   & $1.0282$ & $7.68$ & $2.86\times10^{-1}$ & $6.54\times10^{-2}$ \\
$10^{9}$   & $1.0556$ & $7.73$ & $1.289$ & $3.02\times10^{-1}$ \\
$10^{10}$  & $1.1286$ & $7.84$ & $2.514$ & $6.05\times10^{-1}$ \\
$10^{11}$  & $1.4354$ & $7.84$ & $2.538$ & $6.18\times10^{-1}$ \\
\midrule
diffusion off & --- & $7.76$ & $4.285$ & $1.109$ \\
measured (C26) & --- & --- & $1.24^{+0.22}_{-0.23}$ & $1.108 \pm 0.030$ \\
\bottomrule
\end{tabular}
\end{table}

Read as a measurement of $K_{zz}$, Table~\ref{tab:scan} says four things.

\begin{itemize}
\item \emph{The response is enormous and monotonic below the plateau.} The
      equivalent width rises by four orders of magnitude between
      $K_{zz} = 0$ and $10^{10}$. In that range the line is a sharp probe.
\item \emph{It saturates.} Between $10^{10}$ and $10^{11}$ the equivalent
      width moves by $2\%$. Above the saturation point the line carries no
      information about $K_{zz}$ at all, because the composition in the
      line-forming region has stopped responding
      (Sect.~\ref{sec:outer}).
\item \emph{It is degenerate with composition.} The equivalent width
      measures the metastable column, which is set by the product of the
      reservoir He/H and the fraction of it that survives the homopause.
      One observable cannot separate two unknowns. The scan was run at
      He/H\,$=0.55$, a value that was itself fitted in the well-mixed limit,
      so the scan measures a slice through the degeneracy rather than the
      degeneracy itself.
\item \emph{The wind barely notices.} $\log_{10}\dot M$ moves by $0.23$\,dex
      over eleven decades of $K_{zz}$, and not monotonically. The eddy term
      is a composition control, not a mass-loss control, so the mass-loss
      rate cannot be used as an independent lever on it.
\end{itemize}

Breaking the degeneracy needs a second observable with a different
dependence on the two unknowns. The natural candidates are a hydrogen line
formed in the same region --- the ratio of a He and an H diagnostic depends
on composition and on separation differently --- or a lower-atmosphere
composition prior, which is what a Phase-E \texttt{base.inp} handoff would
supply (Sect.~\ref{sec:next}). Neither exists for this planet at present.
As things stand the He\,$10830$ line constrains the \emph{product} of
reservoir composition and eddy mixing, and quoting a $K_{zz}$ from it alone
would be circular.

\section{Application: the LHS\,1140\,b wind base}
\label{sec:base}

\subsection{The base, and what the eddy term competes against}

The run this is for has its base at $p = 1.0$\,dyn\,cm$^{-2}
= 1.0\,\mu$bar, $T_0 = 226$\,K, derived $n_0 = 3.2049\times10^{13}$\,cm$^{-3}$,
on a planet with $R_p = 1.730\,R_\oplus$, $M_p = 5.60\,M_\oplus$ and
$g = 1837$\,cm\,s$^{-2}$ (\path{LHS1140b/kzz_literature.md}, header).
The Banks \& Kockarts hard-sphere binary coefficient the operator itself
uses,
%
\begin{equation}
  D_{12} = 1.52\times10^{18}\,(1/A_s + 1/A_t)^{1/2}\,T^{1/2}/n ,
\end{equation}
%
evaluates at those conditions to $D({\rm H,He}) =
8.0\times10^{5}$\,cm$^{2}$\,s$^{-1}$, and the run's own stage-resolved
$D_{\rm eff}$ at the first face agrees, $4.6\times10^{5}$.

The crucial structure is not that number but its gradient. From
\texttt{kzz\_decision.md}, section~1, $D_{\rm eff}$ rises from
$4.6\times10^{5}$ at the base to $2.5\times10^{10}$\,cm$^{2}$\,s$^{-1}$ at
$1.2\,R_p$ and $1.4\times10^{11}$ at $2\,R_p$ --- nearly five decades over
one fifth of a planetary radius, driven by the falling density (a factor
$500$ in pressure over that interval) and the rising temperature. Against a
\emph{constant} $K_{zz}$ this has an immediate consequence:
%
\begin{quote}
any $K_{zz}$ above $\sim 10^{6}$ mixes the base cell, and no $K_{zz}$ below
$\sim 10^{11}$ mixes anything at $2\,R_p$.
\end{quote}
%
The choice of \texttt{He\_Kzz} therefore sets \emph{how far above the base}
the homopause sits and nothing else; it never decides whether the base
separates.

\subsection{The literature at $1\,\mu$bar}

Table~\ref{tab:lit} collects the published values evaluated or extrapolated
to the base pressure of this run, with the verification status carried
across from \texttt{kzz\_literature.md}.

\begin{table}[htbp]\centering
\caption{Published $K_{zz}$ evaluated at $1\,\mu$bar. ``Verified'' means the
quoted wording and numerals were read from the publisher's PDF;
``ADS full text'' means they were read from the arXiv version with key
sentences confirmed present in the ADS full-text index of the published
article. Full detail, including which sentences, is in
\protect\path{LHS1140b/kzz_literature.md}.}
\label{tab:lit}
\small
\begin{tabular}{llrll}
\toprule
source & object & $K_{zz}(1\,\mu$bar$)$ & inside stated range? & verified? \\
 & & [cm$^{2}$\,s$^{-1}$] & & \\
\midrule
Parmentier et al.\ (2013) & HD\,209458\,b & $5\times10^{11}$ &
  yes & ADS full text \\
Moses et al.\ (2011) & hot Jupiters & $10^{11}$--$4\times10^{11}$ &
  yes, by assumption & ADS full text \\
Charnay et al.\ (2015) & GJ\,1214\,b & $1.8$--$7.5\times10^{9}$ &
  no ($1.5$ dec.\ above top) & verified \\
Taylor et al.\ (2025) & HD\,209458\,b & $10^{9}$ &
  yes (quoted at $10^{-6}$\,bar) & verified \\
Arfaux \& Lavvas (2023) & \emph{LHS\,1140\,b} & $2.8\times10^{7}$--$1.4\times10^{8}$ &
  yes (plateau) & ADS full text \\
Blain et al.\ (2021) & K2-18\,b & $10^{6}$ nom., $10^{5}$--$10^{10}$ &
  not pressure-resolved & ADS full text \\
\midrule
Festou et al.\ (1981) & \emph{Jupiter, measured} & $1.4\times10^{6}$ &
  at the homopause & abstract (Sect.~\ref{sec:verif}) \\
\bottomrule
\end{tabular}
\end{table}

The spread is five decades, and it is not random: it sorts by how hot and
how strongly forced the source atmosphere is. The two entries that are not
hot-Jupiter extrapolations --- the scaling evaluated for this planet, and
the one measurement in the table --- are also the two lowest. Cherubim et
al.\ (2026), the paper the LHS\,1140\,b measurement comes from, contains no
eddy term, no molecular diffusion and no homopause anywhere in its text or
supplement, so it offers no constraint of its own; its H:He is a fitted
constant of a homogeneous Parker wind.

\subsection{The outer wind, where none of this applies}
\label{sec:outer}

One result of the scan deserves to be separated from the choice of value,
because it is a statement about the physics rather than about the number.
Even at $K_{zz} = 10^{11}$\,cm$^{2}$\,s$^{-1}$ --- above every published
value in Table~\ref{tab:lit} --- the elemental ratio $({\rm He/H})/{\rm
HeH}$ freezes at $0.43$ above $5\,R_p$, and at the adopted $10^{9}$ it
freezes at $0.16$. Since half the $n(2\,^3S)\,r\,dr$ integral lies above
$4.2\,R_p$, most of the He\,$10830$ line is made in gas that no eddy
coefficient reaches. That is why Table~\ref{tab:scan} saturates.

The reason is visible in the $D_{\rm eff}$ profile: molecular diffusion
overtakes any constant eddy term within a few hundredths of a radius of the
base and stays four to five decades ahead of it thereafter. But there is a
second reason, and it is conceptual rather than numerical. Above the
homopause the atmosphere is not a turbulent fluid at all; it is a
transonic outflow in which composition is set by the competition between
advection and molecular separation. Carrying a large constant $K_{zz}$ out
to $30\,R_p$ imposes turbulent mixing in a region where the closure of
Eq.~(\ref{eq:def}) has no eddy field to describe. The operator's behavior
there is dominated by molecular diffusion, so the imposition is currently
harmless; but it means a large \texttt{He\_Kzz} should be read as a
statement about the base region only, and should never be tuned to fix
outer-wind composition.

\subsection{The adopted value, and the weakness of its basis}

\path{LHS1140b/kzz_decision.md} adopts $\texttt{He\_Kzz} =
1.0\times10^{9}$\,cm$^{2}$\,s$^{-1}$, on four grounds: it is the only value
with a published justification for a model of the same kind at the same
pressure (Taylor et al.\ 2025, verified against the published PDF); it is
what EXHALE already uses in \path{examples/14_diffusion} and the
HD\,209458\,b Phase-D baselines; it sits between the two literature
clusters; and it sits on the steep part of the response rather than the
plateau, so its influence is visible to a reader.

Every one of those is an argument from convention or from convenience, and
none is an argument from LHS\,1140\,b. Stated plainly:

\begin{itemize}
\item The one estimate derived from this planet's own gravity and
      temperature is $2.8\times10^{7}$--$1.4\times10^{8}$, one to two
      decades lower (Sect.~\ref{sec:homo}).
\item If its Jovian anchor were taken from the measurement rather than from
      the modelling value, it would be lower again by a factor $\simeq 7$
      (Sect.~\ref{sec:homo-obs}).
\item All three theoretical branches --- convective, wave-driven and
      homopause --- predict that a colder, more weakly forced, higher-gravity
      atmosphere has a smaller $K_{zz}$ (Sect.~\ref{sec:theory}). Nothing in
      the surveyed literature argues the other way.
\item The observable moves by a factor $\simeq 9$ in equivalent width
      between $10^{8}$ and $10^{10}$ (Table~\ref{tab:scan}).
\end{itemize}

The honest form for a paper is therefore the adopted value \emph{and} the
$10^{8}$--$10^{10}$ bracket with the line strengths it gives, which is what
\texttt{kzz\_decision.md} section~5 recommends.

\subsection{Three steps toward a determination}
\label{sec:next}

In increasing order of cost and of credibility:

\paragraph{1. Report the Arfaux--Lavvas axis alongside the adopted value.}
No new physics and no new runs beyond the existing scan: quote
$\texttt{He\_Kzz} = 10^{9}$ as the hot-Jupiter convention and
$10^{8}$ as the planet-specific homopause scaling, with both line
strengths. This costs nothing and removes the main criticism of the present
choice, which is not that the value is wrong but that its provenance is
invisible.

\paragraph{2. Give $K_{zz}$ a profile instead of a constant --- done for the
handoff, and measured.} The physics of Sect.~\ref{sec:gw} does not produce a
constant; it produces growth as
%
\begin{equation}
  K_{zz}(P) = K_{\rm ref} \left(\frac{P}{P_{\rm ref}}\right)^{-\alpha},
  \qquad \alpha = 1/2 ,
  \label{eq:kzzpow}
\end{equation}
%
up to a saturation level, and a turnover above it.

A radially varying $K_{zz}$ has not been hypothetical in the code since
Phase~E: the operator reads a cell array \texttt{kzz\_cell}
(\path{src/modules/init/parameters.f90}), which a
\texttt{Lower atmosphere profile:} file fills by interpolating its own
$K_{zz}(p)$ column onto the grid
(\path{src/modules/files_IO/lower_atmosphere_profile.f90}) and which is
filled uniformly from \texttt{He\_Kzz} otherwise. What was missing was a way
to \emph{state} Eq.~(\ref{eq:kzzpow}); that is now there, on the handoff
side. \texttt{eddy\_diffusion\_coefficient}
(\path{src/utils/lower_profile_schema.py}) takes \texttt{-{}-kzz-power}
$\alpha$ with its amplitude \texttt{-{}-kzz-ref} $K_{\rm ref}$ and the
pressure that amplitude is stated at, \texttt{-{}-kzz-ref-bar}
$P_{\rm ref}$, and writes the resulting column into the profile. In
\texttt{photochem\_to\_lower\_profile.py} the stated coefficient is the one
the chemistry is \emph{solved} on and not only the one the file reports, so
the composition handed over is the composition that $K_{zz}(p)$ produces; in
the VULCAN cross-check arm the chemistry is already finished, so a stated
coefficient changes only what is handed over. With neither option given
each adapter keeps whatever $K_{zz}(p)$ its own solution carried, and the
\texttt{solution\_id} of a constant-$K_{zz}$ run is unchanged.

\paragraph{What the pressure form was then measured to be worth.} On
LHS\,1140\,b, at the two reservoirs of the flux-closed ladder that straddle
the measurement, Eq.~(\ref{eq:kzzpow}) with $\alpha = 1/2$ and
$K_{\rm ref} = 10^{9}$ anchored at the top of the profile
($1.1\times10^{-8}$\,bar) --- which leaves the value the wind inherits above
the file unchanged and gives $1.05\times10^{8}$ at the $1\,\mu$bar match,
$1.05\times10^{5}$ at $1$\,bar and $2.56\times10^{4}$ at the $16.73$\,bar
deep boundary --- moves the He\,10830 equivalent width by $+0.28\%$ when only
the chemistry sees it and $+0.10\%$ when the wind sees it too
(He/H\,$=11.11$ arm), and by $+0.61\%$ and $-0.15\%$ on the He/H\,$=2.09$
arm. The FWHM and $\log_{10}\dot M$ do not move at all. The reason is
visible one step earlier: the elemental ratios at the match, which with the
temperature are what EXHALE consumes from the profile, agree with the
constant-$K_{zz}$ ones to $1.2\%$ or better ($X_{\rm He}$ $0.9985$,
$X_{\rm C}$ $0.990$, $X_{\rm N}$ $0.988$, $X_{\rm O}$ $0.982$--$0.989$),
even though the \emph{carriers} move by up to four decades and nitrogen
changes carrier outright from NH$_3$ to N$_2$. Re-solving the same profile
from a different seed moves the equivalent width by $1.5\%$, an order of
magnitude more than the eddy form does. Measurement and runs:
\path{LHS1140b/exhale/kzz_power/}, recorded in
\path{LHS1140b/kzz_decision.md} section~9 and
\path{docs/lhs1140b_exhale_vs_pwinds.tex}, ``What the constant $K_{zz}$
carries''.

\paragraph{The real remaining gap is the two-branch law, not the profile.}
Three things are still absent, in decreasing order of how much they matter.

\emph{(i) There is no convective branch.} Equation~(\ref{eq:kzzpow}) is the
\emph{stratified} law, and on LHS\,1140\,b the climate solve puts the
tropopause at $2.636$\,bar; below it the mixing-length form of
Eq.~(\ref{eq:gc}) applies and $K_{zz}$ should flatten or grow with depth
rather than collapse. Running the single power law to $16.73$\,bar gives
$1.05\times10^{5}$\,cm$^2$\,s$^{-1}$ at $1$\,bar, at or below the floor of
the modelling literature --- one and a half to two and a half decades under
the $3\times10^{6}$--$3\times10^{7}$ of Charnay et al.\ (2015), $3.7$
decades under the $5\times10^{8}$ of Parmentier et al.\ (2013), and equal to
the $10^{5}$ minimum Ackerman \& Marley (2001) impose by hand. What is
needed is the pair: a wave law above the radiative--convective boundary
joined to a mixing-length law below it, with the boundary taken from the
climate solve that is already in the same adapter.

\emph{(ii) There is no ceiling.} An unbounded $P^{-1/2}$ carried far above
the level its amplitude was set at reaches absurd values and contradicts
Lindzen's own turnover result. The option is deliberately confined to the
handoff, where the file's own top bounds it and where EXHALE gives every
cell above that top the top level's value. That bound is also what limits
the measurement above on the wind side: the LHS\,1140\,b column reaches only
$1.0091\,R_p$ past the match, so the coefficient the wind saw was changed in
a shell shallower than the $1.056\,R_p$ homopause and not where helium
settles. A caller who anchors the same law at the match instead of at the
top raises the coefficient over the whole wind domain by the same factor and
should read that as a change of magnitude, not of form.

\emph{(iii) A run with no profile file still has only the constant}
\texttt{He\_Kzz}. Nothing about Eq.~(\ref{eq:kzzpow}) has been put into
\texttt{binary\_element\_diffusion.f90} itself.

None of this repairs Sect.~\ref{sec:outer}: $D_{\rm eff}$ rises faster than
$P^{-1/2}$ in this wind (roughly as $T^{1/2}/n$, and $T$ rises steeply
through the same interval), so no pressure form of the eddy coefficient
reaches the region where most of the line is made. The gain is not a
stronger line; it is that the base region is described by the physics the
literature values were derived from. And the measurement above says what the
gain is worth on the observable, which is very little --- while noting that
the anchor's own value at the match, $1.05\times10^{8}$, is better founded
for this planet than the adopted $10^{9}$, since it falls inside the
$2.8\times10^{7}$--$1.4\times10^{8}$ of Eq.~(\ref{eq:al}).

\paragraph{3. A determination, by GCM or by simultaneous fit.} Either
measure $K_{zz}$ for this planet the way Sect.~\ref{sec:gcm} describes ---
a GCM of a temperate, helium-rich, tidally locked super-Earth with passive
tracers, which does not currently exist --- or fit $K_{zz}$ jointly with the
reservoir composition against the measured line, accepting that a single
observable cannot separate the two (Sect.~\ref{sec:he-constraint}) and that
the fit would therefore need a second constraint. The Phase-E route is a
partial version of the second, and it is now built: a
\texttt{Lower atmosphere profile:} file carries the composition
\emph{and} $K_{zz}(p)$ from the photochemical model, so both reach the wind
from one calculation instead of two hand-set numbers. The scalar
\texttt{Kzz\_base} key of \texttt{read\_base\_inp} remains for a run handed
over by scalars alone, and is refused when a profile is in use
(\texttt{refuse\_scalar\_key},
\path{src/modules/files_IO/input_read.f90}).

\medskip
\noindent None of the three makes $K_{zz}$ a derived quantity of EXHALE.
It stands for turbulence and wave breaking that a one-dimensional wind model
does not resolve, so it is stated, not computed --- which is why the code's
default is $0$ rather than an inherited number.

\section{Verification status of the new citations}
\label{sec:verif}

The citations below were not in \path{LHS1140b/kzz_literature.md} and
were checked for this memo. Publisher PDFs could not be retrieved for three
of the four: the ADS link gateway returned HTTP~404 for
\texttt{PUB\_PDF} on the two \emph{JGR} papers and on the \emph{Science}
paper, and HTTP~400 on the \emph{ApJ} paper. Where the PDF was
unavailable, quoted wording was checked by (i) the abstract of the published
article as indexed by ADS and (ii) \texttt{full:"\ldots"} phrase queries
against the ADS full-text index of the published article, which is the same
procedure used in \texttt{kzz\_literature.md}. All four papers are
full-text indexed.

\begin{table}[htbp]\centering
\caption{New citations and how each was checked.}
\label{tab:verif}
\begin{tabular}{llll}
\toprule
reference & bibcode & edition read & status \\
\midrule
Lindzen (1981) & \texttt{1981JGR....86.9707L} & published abstract
  & partial \\
Ackerman \& Marley (2001) & \texttt{2001ApJ...556..872A} &
  arXiv:astro-ph/0103423 & partial \\
Prinn \& Barshay (1977) & \texttt{1977Sci...198.1031P} &
  published abstract & partial \\
Festou et al.\ (1981) & \texttt{1981JGR....86.5715F} &
  published abstract & partial \\
\bottomrule
\end{tabular}
\end{table}

\paragraph{Lindzen (1981).} Every sentence quoted in Sect.~\ref{sec:gw} is
from the abstract of the published article as indexed by ADS. Two of them
were additionally confirmed present in the ADS full-text index of the
published article: \texttt{full:"turbulence in the upper mesosphere arises
from the unstable breakdown of tides and gravity waves"} and
\texttt{full:"prevent the growth of wave amplitude with height"} each return
the paper; so does \texttt{full:"eddy coefficients may peak"}.
\textbf{Caveat:} the full-text index of this paper appears to be partial or
OCR-degraded --- \texttt{full:"eddy diffusion"}, \texttt{full:"Richardson
number"} and \texttt{full:"vertical wavelength"} all return nothing, while
\texttt{full:"eddy coefficient"}, \texttt{full:"critical level"} and
\texttt{full:"momentum deposition"} return the paper --- so absence of a
phrase is not evidence here. Lindzen's saturation formula for the eddy
coefficient is deliberately \emph{not} written out in this memo: it was
\textbf{not verified against the published text}, and the physical argument
in Sect.~\ref{sec:gw} is stated in words that the abstract supports.

\paragraph{Ackerman \& Marley (2001).} Read from the arXiv PDF
(\path{astro-ph/0103423}, retrieved through the ADS
\texttt{EPRINT\_PDF} gateway). Eq.~(\ref{eq:gc}) is their eq.~(5), with the
attribution to Gierasch \& Conrath (1985) in their own text. The following
were confirmed present in the ADS full-text index of the published
\emph{ApJ} article: ``is assumed to be the same as that for heat as derived
for free convection''; ``the constant coefficient scaling the eddy diffusion
coefficient is only loosely constrained by observations''; ``the mixing
length is typically assumed to be the pressure scale height''; ``the
convective velocity scale from mixing-length theory''; ``we also assume an
eddy diffusion coefficient no less than a prescribed minimum value''.
\textbf{Caveat:} the composite phrase ``which represents residual
turbulence due to breaking buoyancy waves'' returns nothing, while
``residual turbulence'' and ``breaking buoyancy waves'' each return the
paper. The most likely cause is that an inline equation splits the sentence
in the published typesetting, but the composite wording is
\textbf{not verified against the published text} and is quoted here only as
its two verified fragments. The numerals ($K_{\rm min} = 10^{5}$,
$H = L = 20$\,km, $K = 2\times10^{8}$\,cm$^{2}$\,s$^{-1}$, $w_* =
1.1$\,m\,s$^{-1}$, $T_{\rm eff} = 124$\,K, prefactor $1/3$) are read from
the arXiv version and are \textbf{not verified against the published text}.

\paragraph{Prinn \& Barshay (1977).} All three quotations in
Sect.~\ref{sec:quench} are from the abstract of the published \emph{Science}
article as indexed by ADS. Two were additionally confirmed present in the
ADS full-text index of the published article: \texttt{full:"sensitive
function of the vertical eddy mixing coefficient"} and \texttt{full:"three
to four orders of magnitude greater than that in the earth"}. No numerical
value of $K_{zz}$ from this paper is quoted in this memo, because none was
verified: \texttt{full:"eddy diffusion coefficient"} and
\texttt{full:"quench"} both return nothing for this bibcode, so the paper
uses different wording and the body could not be read.

\paragraph{Festou et al.\ (1981).} The measured homopause coefficient
$1.4^{+0.8}_{-0.7}\times10^{6}$\,cm$^{2}$\,s$^{-1}$, the equatorial
qualifier, the consistency with the H\,Ly$\alpha$ and He\,$584$\,\AA{}
determination, the occultation date, the star, and the wavelength ranges are
all from the abstract of the published article as indexed by ADS. The phrase
\texttt{full:"yields a value of the eddy diffusion coefficient at the
homopause"} returns the paper in the ADS full-text index of the published
article, so the abstract wording is present in the article itself.
The body of the paper was not read.

\paragraph{Not verified.} Two further attributions are carried over from
\texttt{kzz\_literature.md} without independent checking here and are marked
in the text: the anchor $K_{zz}^{J} = 10^{7}$\,cm$^{2}$\,s$^{-1}$ that
Arfaux \& Lavvas (2023) attribute to Koskinen et al.\ (2010)
(Sect.~\ref{sec:homo-obs}) --- the specific Koskinen et al.\ (2010) paper
was not identified --- and the identification of Eq.~(\ref{eq:gc}) with
eq.~(16) of Charnay et al.\ (2015).

\section{References}

\begin{itemize}\itemsep2pt

\item Ackerman, A.~S. \& Marley, M.~S. 2001, ``Precipitating Condensation
Clouds in Substellar Atmospheres'', \textit{ApJ}, 556, 872.
Bibcode \texttt{2001ApJ...556..872A}; DOI \path{10.1086/321540}.
Read from arXiv:astro-ph/0103423 (Sect.~\ref{sec:verif}).

\item Arfaux, A. \& Lavvas, P. 2023, ``A physically derived eddy
parametrization for giant planet atmospheres with application to
exoplanets'', \textit{MNRAS}, 522, 2525.
Bibcode \texttt{2023MNRAS.522.2525A};
\path{references/Arfaux_2023_MNRAS_522_2525_arXiv.pdf}
(arXiv:2304.06314).

\item Banks, P.~M. \& Kockarts, G. 1973, \textit{Aeronomy} (New York:
Academic Press). The source of the $(D+K)$ transport form of
Eq.~(\ref{eq:DplusK}) and of the hard-sphere binary coefficient
implemented in \texttt{binary\_element\_diffusion.f90}.

\item Blain, D., Charnay, B. \& B\'ezard, B. 2021, ``1D atmospheric study of
the temperate sub-Neptune K2-18b'', \textit{A\&A}, 646, A15.
Bibcode \texttt{2021A\&A...646A..15B};
\path{references/Blain_2021_A+A_646_A15_arXiv.pdf}
(arXiv:2011.10459).

\item Chapman, S. \& Cowling, T.~G. 1970, \textit{The Mathematical Theory of
Non-uniform Gases}, 3rd ed.\ (Cambridge: Cambridge Univ.\ Press).
Their eq.~(18.2,6) is the diffusion equation the EXHALE operator solves in
its binary case.

\item Charnay, B., Meadows, V. \& Leconte, J. 2015, ``3D Modeling of
GJ1214b's Atmosphere: Vertical Mixing Driven by an Anti-Hadley
Circulation'', \textit{ApJ}, 813, 15. Bibcode \texttt{2015ApJ...813...15C};
\path{references/Charnay_2015_ApJ_813_15.pdf} (publisher PDF).

\item Cherubim, C. et al.\ 2026, ``Helium escaping from the atmosphere of a
nearby rocky exoplanet orbiting in a habitable zone'', \textit{Science}
(first release, 16 July 2026).
\path{references/Cherubim_2026Science.pdf} and its supplement
(publisher PDF).

\item Festou, M.~C., Atreya, S.~K., Donahue, T.~M., Sandel, B.~R., Shemansky,
D.~E. \& Broadfoot, A.~L. 1981, ``Composition and Thermal Profiles of the
Jovian Upper Atmosphere Determined by the Voyager Ultraviolet Stellar
Occultation Experiment'', \textit{J.\ Geophys.\ Res.}, 86, 5715.
Bibcode \texttt{1981JGR....86.5715F}; DOI
\path{10.1029/JA086iA07p05715}. (The first four authors are those
returned by ADS; the remainder of the author list was not checked.)

\item Gierasch, P.~J. \& Conrath, B.~J. 1985, ``Energy conversion processes
in the outer planets'', in \textit{Recent Advances in Planetary
Meteorology}, ed.\ G.~E.\ Hunt (Cambridge: Cambridge Univ.\ Press), 121.
Bibcode \texttt{1985rapm.book..121G}. Cited through Ackerman \& Marley
(2001) and Charnay et al.\ (2015); not read.

\item Koskinen, T.~T. et al.\ 2013a, ``The escape of heavy atoms from the
ionosphere of HD209458b. I.'', \textit{Icarus}, 226, 1678.
Bibcode \texttt{2013Icar..226.1678K};
\path{references/Koskinen_2013Icarus_226_1678.pdf}.

\item Koskinen, T.~T. et al.\ 2013b, ``The escape of heavy atoms from the
ionosphere of HD209458b. II.'', \textit{Icarus}, 226, 1695.
Bibcode \texttt{2013Icar..226.1695K};
\path{references/Koskinen_2013Icarus_226_1695.pdf}.

\item Koskinen, T.~T., Lavvas, P., Huang, C. et al.\ 2022, \textit{ApJ},
929, 52. Bibcode \texttt{2022ApJ...929...52K};
\path{references/Koskinen_2022_ApJ_929_52.pdf}. $K_{zz}$ appears only
as a symbol in its Appendix~B diffusion equation
($\Lambda_s = K_{zz}/D_s$); no numerical value appears in the extractable
text.

\item Lindzen, R.~S. 1981, ``Turbulence and Stress Owing to Gravity Wave and
Tidal Breakdown'', \textit{J.\ Geophys.\ Res.}, 86, 9707.
Bibcode \texttt{1981JGR....86.9707L}; DOI
\path{10.1029/JC086iC10p09707}.

\item Moses, J.~I. et al.\ 2011, ``Disequilibrium Carbon, Oxygen, and
Nitrogen Chemistry in the Atmospheres of HD 189733b and HD 209458b'',
\textit{ApJ}, 737, 15. Bibcode \texttt{2011ApJ...737...15M};
\path{references/Moses_2011_ApJ_737_15_arXiv.pdf}
(arXiv:1102.0063).

\item Parmentier, V., Showman, A.~P. \& Lian, Y. 2013, ``3D mixing in hot
Jupiter atmospheres. I. Application to the day/night cold trap'',
\textit{A\&A}, 558, A91. Bibcode \texttt{2013A\&A...558A..91P};
\path{references/Parmentier_2013_A+A_558_A91_arXiv.pdf}
(arXiv:1301.4522).

\item Prinn, R.~G. \& Barshay, S.~S. 1977, ``Carbon Monoxide on Jupiter and
Implications for Atmospheric Convection'', \textit{Science}, 198, 1031.
Bibcode \texttt{1977Sci...198.1031P}; DOI
\path{10.1126/science.198.4321.1031}.

\item Taylor, J. et al.\ 2025, \textit{ApJ}, 989, 68.
Bibcode \texttt{2025ApJ...989...68T};
\path{references/Taylor_2025_ApJ_989_68.pdf} (publisher PDF).

\item Zhang, X. \& Showman, A.~P. 2018a, ``Global-mean Vertical Tracer Mixing
in Planetary Atmospheres. I. Theory and Fast-rotating Planets'',
\textit{ApJ}, 866, 1. Bibcode \texttt{2018ApJ...866....1Z};
\path{references/Zhang_2018_ApJ_866_1_arXiv.pdf}
(arXiv:1803.09149).

\item Zhang, X. \& Showman, A.~P. 2018b, ``Global-mean Vertical Tracer
Mixing in Planetary Atmospheres. II. Tidally Locked Planets'',
\textit{ApJ}, 866, 2. Bibcode \texttt{2018ApJ...866....2Z};
\path{references/Zhang_2018_ApJ_866_2_arXiv.pdf}
(arXiv:1808.05365).

\end{itemize}

\medskip
\noindent\textbf{Companion documents in this tree.}
\path{LHS1140b/kzz_literature.md} (the survey, with the verification
status of every published number);
\path{LHS1140b/kzz_decision.md} (base conditions, the $D_{\rm eff}$
structure, the \texttt{He\_Kzz} scan, the adopted value and the
alternatives); \path{docs/binary_diffusion_design.md} (the operator,
sections~2.3 and~9.4); \path{docs/lhs1140b_exhale_vs_pwinds.tex}
(the He\,$10830$ comparison the composition was fitted against).

\end{document}
